\subsection{形式幂级数环中方程的解}

引理 2. --- 设 \((g_i)_{i \in I}\) 是一族阶 \(\geqslant 2\) 的元素，位于 A{[}{[}I{]}{]} 中。当 I 为无限时，假设 \(g_i\) 沿 I 的有限子集的补集构成的滤子趋于 0 。存在且仅存在拓扑 A-代数 A{[}{[}I{]}{]} 的一个自同构 T ，使得对所有 i ∈ I，有 T(X \(_i\) ) = X \(_i\) + g \(_i\) 。进一步，

\[
\omega (\mathrm{T} (u) - u) \geqslant \omega (u) + 1\tag{23}
\]

对 \(u \in \mathbf{A}[[\mathrm{I}]]\) 成立。

级数 \(f_{i} = \mathbf{X}_{i} + g_{i}\) 没有常数项，且当 I 无限时 \(f_{i}\) 沿 I 的有限子集的补集构成的滤子趋于 0 。因此（IV，第28页，命题 4）存在恰好一个 A-代数 A{[}{[}I{]}{]} 的连续自同态 T ，使得 \(\mathrm{T}(\mathbf{X}_i) = f_i\) 对所有 \(i \in I\) 。对每个 \(\nu \in \mathbf{N}^{(I)}\) 我们令

\[
v _ {\nu} = \mathrm{T} \left(\mathbf {X} ^ {\nu}\right) - \mathbf {X} ^ {\nu} = \prod_ {i \in I} \left(\mathbf {X} _ {i} + g _ {i}\right) ^ {\nu (i)} - \prod_ {i \in I} \mathbf {X} _ {i} ^ {\nu (i)};
\]

关系 \(\omega(g_i) \geqslant 2\) 蕴含 \(\omega(v_\nu) \geqslant |\nu| + 1\) ，并且关系式 (23) 由此立即得出。

让我们证明 T 是单射。给定 \(u \in A[[I]]\) 使得 \(T(u) = 0\) ，由 (23) 我们有 \(\omega(u) \geqslant \omega(u) + 1\) ，如果 \(u \neq 0\) 则这是不可能的，因为 \(\omega(u)\) 将是一个正整数。

对每个形式级数 \(v\) （位于 \(A[[I]]\) 中），我们用 \(H_n(v)\) 表示它的 \(n\) 次齐次分量。令 \(S_0(v) = H_0(v)\) ，并用递推方程定义连续映射 \(S_n: A[[I]] \to A[[I]]\)

\[
\mathrm{S} _ {n} (v) = \mathrm{H} _ {n} \left(v - \mathrm{T} \left(\sum_ {k = 0} ^ {n - 1} \mathrm{S} _ {k} (v)\right)\right) \quad \text { for } \quad n \geqslant 1.\tag{24}
\]

令 \(S(v) = \sum_{n \geqslant 0} S_n(v)\) ；如果 \(\nu \in \mathbf{N}^{(I)}\) 且 \(n = |\nu|\) ，则系数 \(S^{\nu}(v)\) ，即

\(X^{\nu}\) 在 \(S(v)\) 中的系数，等于 \(X^{\nu}\) 在 \(S_{n}(v)\) 中的系数；由于 \(S_{n}\) 是连续映射，映射 \(S^{\nu}: A[[I]] \to A\) 是连续的。因此，根据 \(A[[I]] = A^{N^{(I)}}\) 上乘积拓扑的定义，映射 \(S: A[[I]] \to A[[I]]\) 是连续的。

我们将证明关系式 \(\mathrm{T}(\mathrm{S}(v)) = v\) 对所有 \(v \in A[[I]]\) 成立，这将完成引理的证明。设 \(v \in A[[I]]\) ，\(u_n = S_n(v)\) 且 \(u = S(v)\) 。设 n 为正整数，使得

\[
\omega (v - \mathrm{T} (u)) \geqslant n.\tag{$(25)_{n}$}
\]

我们有 \(\omega(u - (u_{0} + \cdots + u_{n-1})) \geqslant n\) ，从而

\[
\omega (\mathrm{T} (u) - \mathrm{T} (u _ {0} + \dots + u _ {n - 1}) - u _ {n}) \geqslant n + 1\tag{26}
\]

由 (23)，现在递推方程 (24) 表明

\[
u _ {n} = \mathrm{H} _ {n} (v - \mathrm{T} (u _ {0} + \dots + u _ {n - 1}))\tag{27}
\]

由 (26)，形式幂级数 \(v - \mathrm{T}(u)\) 与 \(v - \mathrm{T}(u_{0} + \cdots + u_{n-1}) - u_{n}\) 有相同的 n 次齐次分量，且由 (27) 该分量为零。于是我们有 \(\omega(v - \mathrm{T}(u)) \geqslant n + 1\) ，即 \((25)_{n}\) 蕴含 \((25)_{n+1}\) 。由于 \((25)_{0}\) 显然成立，我们遂有 \(\omega(v - \mathrm{T}(u)) \geqslant n\) 对每个整数 \(n \geqslant 0\) ，从而 \(v = \mathrm{T}(u) = \mathrm{T}(\mathrm{S}(v))\) ，此即所要证明的。

在本节其余部分中，对于任意集合 I，我们将用 A \{I\} 表示满足以下条件的族 \((f_i)_{i\in I}\) 的集合：

\begin{enumerate}
\def\labelenumi{(\roman{enumi})}
\item
  对每个 \(i \in I\) ，\(f_i\) 是 A{[}{[}I{]}{]} 的一个无常数项的元素；
\item
  若 I 是无限的，则 \(f_{i}\) 沿 I 的有限子集的补集所构成的滤子趋于 0。
\end{enumerate}

集合 A\{I\} 对复合律 \((\mathbf{f},\mathbf{g})\mapsto\mathbf{f}\circ\mathbf{g}\) 成一个幺半群，以 \(\{X_{i}\}_{i\in I}\) 为单位元。因此 A\{I\} 的可逆元素之集是一个群。

另一方面，设 E 为 A-代数 \(\mathbf{A}[[\mathrm{I}]]\) 的所有保持无常数项形式幂级数组成的理想不变的连续保单位自同态构成的幺半群。如果 \(\mathbf{f} \in \mathbf{A}\{\mathbf{I}\}\) 且 \(g \in \mathbf{A}[[\mathrm{I}]]\) ，则元素 \(g(\mathbf{f})\) 有定义。对固定的 \(\mathbf{f}\) ，映射 \(g \mapsto g(\mathbf{f})\) （从 \(\mathbf{A}[[\mathrm{I}]]\) 到其自身）是 E 的一个元素 \(\mathbf{W}_{\mathbf{f}}\) 。如果 \(\mathbf{f}_1, \mathbf{f}_2 \in \mathbf{A}\{\mathbf{I}\}\) 且 \(g \in \mathbf{A}[[\mathrm{I}]]\) ，由公式 (5)（IV，第30页）我们有

\[
\mathbf {W} _ {\mathbf {f} _ {1} \circ \mathbf {f} _ {2}} (g) = g (\mathbf {f} _ {1} \circ \mathbf {f} _ {2}) = g (\mathbf {f} _ {1}) \circ \mathbf {f} _ {2} = \mathbf {W} _ {\mathbf {f} _ {2}} (\mathbf {W} _ {\mathbf {f} _ {1}} (g))
\]

因此 \(\mathbf{f} \mapsto \mathbf{W}_{\mathbf{f}}\) 是从与 A \{I\} 相反的幺半群到 E 中的一个同态。由命题 4（IV，第28页），此同态是双射的。

设 \(\mathbf{f} = (f_i)_{i\in I}\in \mathbf{A}\{\mathbf{I}\}\) ，并令 \(\sum_{j\in I}\alpha_{ij}\mathbf{X}_j\) 为

\(f_{i}\) 的 1 次齐次分量。对 I 中任意固定的 \(j\) ，我们有 \(\alpha_{ij} = 0\) ，除有限个下标 \(i\) 外，这由上述假设 (ii) 可知。若 \((\lambda_i) \in \mathbf{A}^{(I)}\) ，我们遂有 \(\left(\sum_{j \in I} \alpha_{ij} \lambda_j\right) \in \mathbf{A}^{(I)}\) 。

我们用 \(T_{f}\) 表示 A-线性映射 \(^{1}\)

\[
\left(\lambda_ {i}\right) \mapsto \left(\sum_ {j \in I} \alpha_ {i j} \lambda_ {j}\right)
\]

（从 \(\mathbf{A}^{(I)}\) 到 \(\mathbf{A}^{(I)}\) ）。若 \(\mathbf{g} \in \mathbf{A}\{I\}\) ，则容易验证

\[
\mathrm{T} _ {\mathbf {f} \circ \mathbf {g}} = \mathrm{T} _ {\mathbf {f}} \circ \mathrm{T} _ {\mathbf {g}}.\tag{28}
\]

命题 10. --- 设 \(\mathbf{f} \in \mathbf{A}\{\mathbf{I}\}\) ；则下列条件等价：(i) \(\mathbf{f}\) 在 \(\mathbf{A}\{\mathbf{I}\}\) 中对运算 \(\circ\) 可逆；

\begin{enumerate}
\def\labelenumi{(\roman{enumi})}
\setcounter{enumi}{1}
\tightlist
\item
  \(\mathbf{T}_{\mathbf{f}}\) 在环 \(\operatorname{End}(\mathbf{A}^{(I)})\) 中可逆。
\end{enumerate}

中可逆。蕴含关系 (i) \(\Rightarrow\) (ii) 由 (28) 直接得出。现在设 \(T_{f}\) 在 \(\operatorname{End}(A^{(I)})\) 中可逆。存在 \(\mathbf{g} = (g_i)_{i \in I} \in A\{I\}\) ，使得每个 \(g_i\) 都是 1 次齐次的，且 \(T_{g} \circ T_{f}\) 是 \(A^{(I)}\) 的恒等映射。记 \(\mathbf{h} = \mathbf{g} \circ \mathbf{f}\) ；则 (28) 表明 \(T_{h}\) 是 \(A^{(I)}\) 的恒等映射，这等价于断言 \(\omega(h_i - X_i) \geqslant 2\) 。由 IV 第35页的引理 2，h 因此在 A \{I\} 中可逆。显然 \(\mathbf{g}\) 在 A \{I\} 中可逆，因此 \(\mathbf{f}\) 在 A \{I\} 中可逆。

推论. --- 设 \(f_{i}(Y_{1}, Y_{2}, \ldots, Y_{q}, X_{1}, X_{2}, \ldots, X_{p}) (1 \leqslant i \leqslant q)\) 是 A{[}{[} \(Y_{1}, \ldots, Y_{q}, X_{1}, \ldots, X_{p}\) {]}{]} 中 q 个无常数项的形式幂级数。若形式幂级数 D = det \(\left(\frac{\partial f_{i}}{\partial Y_{j}}\right)\) 的常数项在 A 中可逆，则存在恰好一组 q 个形式幂级数 \(u_{1}(X_{1}, \ldots, X_{p}), \ldots, u_{q}(X_{1}, \ldots, X_{p})\) ，使得

\[
f _ {i} (u _ {1}, \dots , u _ {q}, \mathrm{X} _ {1}, \dots , \mathrm{X} _ {p}) = 0 \quad (1 \leqslant i \leqslant q).\tag{29}
\]

令 \(f_{q+1} = \mathbf{X}_1, \ldots, f_{q+p} = \mathbf{X}_p\) ，\(\mathbf{f} = (f_1, \ldots, f_{p+q})\) ，则 \(\det \mathbf{T}_{\mathbf{f}}\) 等于 \(\mathbf{D}\) 的常数项，从而在 \(\mathbf{A}\) 中可逆；于是 \(\mathbf{T}_{\mathbf{f}}\) 可逆。由命题 10，存在无常数项的形式幂级数

\[
\boldsymbol {g} _ {1}, \dots , \boldsymbol {g} _ {q + p} \in \mathrm{A} \left[ \left[ \mathrm{Y} _ {1}, \dots , \mathrm{Y} _ {q}, \mathrm{X} _ {1}, \dots , \mathrm{X} _ {p} \right] \right]
\]

使得当记

\[
\mathbf {g} = \left(g _ {1}, \dots , g _ {p + q}\right), \quad \mathbf {1} _ {p + q} = \left(\mathrm{Y} _ {1}, \dots , \mathrm{Y} _ {q}, \mathrm{X} _ {1}, \dots , \mathrm{X} _ {p}\right)
\]

时，我们有 \(\mathbf{f} \circ \mathbf{g} = \mathbf{g} \circ \mathbf{f} = \mathbf{1}_{p + q}\) 。关系式 \(\mathbf{f} \circ \mathbf{g} = \mathbf{1}_{p + q}\) 特别地给出

\[
\boldsymbol {g} _ {q + 1} = \mathbf {X} _ {1}, \dots , \boldsymbol {g} _ {q + p} = \mathbf {X} _ {p}.
\]

因此

\[
f _ {i} \left(g _ {1}, \dots , g _ {q}, \mathrm{X} _ {1}, \dots , \mathrm{X} _ {p}\right) = \mathrm{Y} _ {i} \quad (1 \leqslant i \leqslant q).\tag{30}
\]

现在令

\[
u _ {i} \left(\mathrm{X} _ {1}, \dots , \mathrm{X} _ {p}\right) = g _ {i} (0, \dots , 0, \mathrm{X} _ {1}, \dots , \mathrm{X} _ {p}) \quad (1 \leqslant i \leqslant q);\tag{31}
\]

以 0 代替每个 \(\mathbf{Y}_i\) 代入 (30) ，我们得到所要的关系式 (29)。

反之，假设形式幂级数 \(u_1, \ldots, u_q\) （属于环 \(A[[X_1, \ldots, X_p]]\) ）满足关系式 (29)。关系式 \(\mathbf{g} \circ \mathbf{f} = 1_{p + q}\) 蕴含

\[
g _ {i} (f _ {1}, \dots , f _ {q}, \mathrm{X} _ {1}, \dots , \mathrm{X} _ {p}) = \mathrm{Y} _ {i} \quad (1 \leqslant i \leqslant q);\tag{32}
\]

并在 (32) 中以 \(u_{i}\) 代 \(Y_{i}\) （对 \(1 \leqslant i \leqslant q\) ），我们得到 (31)，由此得出方程组 (29) 的解的唯一性。

